\newglossaryentry{architectureascode}{
    name={Architecture-as-Code},
    description={подход к представлению архитектуры программной системы в машинно обрабатываемом виде, пригодном для версионирования, анализа и использования в автоматизированных процессах}
}

\newglossaryentry{infrastructureascode}{
    name={Infrastructure-as-Code},
    description={подход к описанию инфраструктурных ресурсов и конфигураций в виде кода или декларативных спецификаций, которые могут храниться в системе контроля версий; такие спецификации применяются автоматизированными инструментами}
}

\newglossaryentry{desiredstate}{
    name={Desired state},
    description={желаемое состояние инфраструктуры, описывающее, какие хосты и workload должны находиться под управлением платформы и какими свойствами они должны обладать}
}

\newglossaryentry{actualstate}{
    name={Actual state},
    description={фактическое состояние инфраструктуры, полученное из источников инвентаризации и наблюдения на момент проверки}
}

\newglossaryentry{drift}{
    name={Drift},
    description={расхождение между desired state и actual state}
}

\newglossaryentry{reconcile}{
    name={Reconcile},
    description={процесс приведения фактического состояния инфраструктуры к желаемому состоянию}
}

\newglossaryentry{controlloop}{
    name={Control loop},
    description={повторяемый цикл управления, включающий получение desired state, получение actual state, сравнение состояний, обнаружение drift и запуск reconcile-действий}
}

\newglossaryentry{inventory}{
    name={Inventory},
    description={представление сведений об управляемых хостах, их атрибутах, доступности источников данных и фактически наблюдаемых workload}
}

\newglossaryentry{workload}{
    name={Workload},
    description={программный компонент или служебный процесс, размещаемый на управляемом хосте}
}

\newglossaryentry{managedhost}{
    name={Managed host},
    description={хост, находящийся в зоне управления платформы и участвующий в цикле сопоставления desired state и actual state}
}

\newglossaryentry{sourceoftruth}{
    name={Source of truth},
    description={источник данных, рассматриваемый системой как авторитетный для определённого типа сведений}
}

\newglossaryentry{partialresult}{
    name={Partial result},
    description={частичный результат получения или обработки данных, при котором часть сведений недоступна}
}

\newglossaryentry{detector}{
    name={Detector},
    description={компонент логики DriftDetectorSvc, отвечающий за проверку конкретного workload и формирование результата обнаружения drift}
}

\newglossaryentry{operator}{
    name={Operator},
    description={компонент логики ReconcilerSvc, отвечающий за построение плана выполнения reconcile для конкретного поддерживаемого workload}
}

\newglossaryentry{executionplan}{
    name={Execution plan},
    description={структурированное описание действия reconcile, включающее выбранный playbook, целевой хост, компонент и параметры выполнения}
}

\newglossaryentry{snapshot}{
    name={Snapshot},
    description={согласованная копия состояния, хранящаяся сервисом в памяти и используемая для выдачи API-ответов без обращения к источникам данных при каждом запросе}
}

\newglossaryentry{structurizr}{
    name={Structurizr},
    description={инструмент и формат описания архитектуры программных систем, поддерживающий C4-модель и экспорт архитектурной модели в JSON-представление}
}

\newglossaryentry{cadvisor}{
    name={cAdvisor},
    description={инструмент сбора сведений о контейнерах и их состоянии, используемый в текущем MVP как источник actual state для контейнерных workload}
}

\newglossaryentry{ansiblerunner}{
    name={Ansible Runner},
    description={инструмент программного запуска Ansible playbook, используемый в ReconcilerSvc для выполнения reconcile-действий}
}

\newglossaryentry{parsersvc}{
    name={ParserSvc},
    description={сервис платформы, отвечающий за разбор Structurizr JSON, построение host-centric desired state и выдачу этого состояния через API}
}

\newglossaryentry{inventorysvc}{
    name={InventorySvc},
    description={сервис платформы, отвечающий за формирование actual state управляемых хостов на основе bootstrap-описания и данных внешних источников}
}

\newglossaryentry{driftdetectorsvc}{
    name={DriftDetectorSvc},
    description={сервис платформы, сравнивающий desired state и actual state, обнаруживающий drift и отправляющий reconcile-команды}
}

\newglossaryentry{reconcilersvc}{
    name={ReconcilerSvc},
    description={сервис платформы, принимающий reconcile-команды, выбирающий operator и асинхронно запускающий действия восстановления}
}
